% ID: SEL-D0001-OA-E03
% ESTADO: cotejada. Fuente y transcripción revisadas de forma independiente; salida compilada inspeccionada.
% Fuente: Selectividad/sel_2019_matematicas_aplicadas/Titular Junio/Exámenes Mat Apli CCSS A y B.pdf
% SHA-256: fa0cbbddfdf34b1f9fcf8d1090c2d0c3f13e5353e06da22ec354d3345e313584
% Página: 1; opción A; ejercicio 3.
% Encabezado original: EJERCICIO 3. Solo enunciado; no contiene solución.

El 65\,\% de los turistas que visitan una provincia elige alojamientos en la
capital y el resto en zonas rurales. Además, el 75\,\% de los turistas que
se hospedan en la capital y el 15\,\% de los que se hospedan en zonas rurales,
lo hacen en hoteles, mientras que el resto lo hace en apartamentos turísticos.
Se elige al azar un turista de los que se han alojado en esa provincia.

\begin{enumerate}[label=\alph*)]
  \item \textbf{(1.5 puntos)} \textquestiondown{}Cuál es la probabilidad de que
  se haya hospedado en un hotel?
  \item \textbf{(1 punto)} Si se sabe que se ha hospedado en un apartamento
  turístico, \textquestiondown{}cuál es la probabilidad de que el apartamento
  esté en zonas rurales?
\end{enumerate}
